\chapter{Anal Pain}
\index{Anal Pain}

\section*{Definition and Overview}
Anal pain is discomfort, aching, burning, or sharp pain felt at the anus or in the surrounding area. It is a symptom rather than a disease, and the underlying causes range from minor and self-limiting problems, such as a small tear or a spasm of the pelvic floor muscles, to more serious conditions that require treatment.

Anal pain is very common and most people will experience it at some point in their lives. Because the area is sensitive and problems can be distressing to discuss, many people delay seeking help; however, most causes are benign and readily treatable, and early assessment is the best way to reach a diagnosis and rule out anything serious.

\section*{Epidemiology}
Anal pain is extremely common, although precise figures are difficult to establish because many people never seek medical advice. Anal fissures and haemorrhoids, two of the most frequent causes, affect people of all ages but are more common in adults. Fissures occur roughly equally in men and women, while haemorrhoids become more frequent with age and during pregnancy.

Constipation and hard stools are major contributing factors in many cases. Painful anal conditions are among the most common reasons for referral to colorectal clinics and for presentation to emergency departments.

\section*{Aetiology and Pathophysiology}
\begin{itemize}
    \item \textbf{Anal fissure:} a tear in the lining of the anal canal, usually from passing a hard or large stool, which goes into spasm and causes sharp pain
    \item \textbf{Haemorrhoids:} swollen blood vessels around the anus or lower rectum that can become inflamed, thrombosed, or prolapsed
    \item \textbf{Muscle spasm:} pelvic floor or sphincter spasm (functional anal pain) occurring with no visible abnormality
    \item \textbf{Perianal abscess or fistula:} infection in the tissues around the anus causing throbbing pain and swelling
    \item \textbf{Infection:} skin infections, herpes, or other sexually transmitted infections causing ulceration
    \item \textbf{Other causes:} anal injury, retained foreign bodies, childbirth tears, inflammatory bowel disease, and rarely anal cancer
    \item \textbf{Mechanism:} the anal canal is richly supplied with nerve endings, so tearing, stretching, spasm, or inflammation in this area is felt as intense pain, often worsened by bowel movements and sitting
\end{itemize}

\section*{Clinical Features}
\begin{itemize}
    \item Sharp or burning pain that may be worse during or after passing stool
    \item Pain that persists or throbs between bowel movements, which is more suggestive of infection
    \item Bleeding, with bright red blood on the toilet paper or in the toilet bowl
    \item A tender lump or swelling at the anal margin
    \item Itching, discharge, or moisture around the anus
    \item Difficulty sitting comfortably
    \item Sometimes a sensation of fullness or an urge to strain
\end{itemize}

\section*{Differential Diagnosis}
\begin{itemize}
    \item Anal fissure
    \item Haemorrhoids (external, internal, or thrombosed)
    \item Perianal abscess or anal fistula
    \item Proctalgia fugax, brief episodes of severe anal pain with no cause found
    \item Levator ani syndrome, or functional pelvic floor pain
    \item Sexually transmitted infections causing anal ulcers, such as herpes or syphilis
    \item Inflammatory bowel disease affecting the anal area, such as Crohn's disease
    \item Anal or rectal cancer, which is uncommon but serious
    \item Referred pain from the back, pelvis, or prostate
\end{itemize}

\section*{When to Seek Medical Care}
See a doctor if anal pain is severe, persists for more than a few days despite simple measures, keeps recurring, or is accompanied by bleeding, a lump, or discharge.

\textbf{Contact your local emergency services for:}
\begin{itemize}
    \item Severe, uncontrollable bleeding from the anus
    \item Sudden severe pain with signs of shock, including faintness, paleness, a rapid heart rate, or collapse
    \item An object that is stuck or deeply embedded in the anal area
\end{itemize}

\section*{Diagnosis and Investigations}
\begin{itemize}
    \item \textbf{History:} how long the pain has been present, what makes it better or worse, any bleeding, and the bowel habit
    \item \textbf{Examination:} inspection of the anal area for tears, lumps, swelling, or skin changes
    \item \textbf{Digital rectal examination:} to assess sphincter tone and feel for masses or tenderness
    \item \textbf{Anoscopy or proctoscopy:} a small instrument used to view the anal canal and lower rectum directly
    \item \textbf{Flexible sigmoidoscopy or colonoscopy:} if bleeding is persistent or inflammatory bowel disease is suspected
    \item \textbf{Imaging:} ultrasound or MRI for suspected deep abscesses or fistulas
    \item \textbf{Swabs or biopsy:} to look for infection or to investigate suspicious lesions
\end{itemize}

\section*{Management}
\begin{itemize}
    \item \textbf{Self-care:} warm baths (sitz baths), gentle cleaning, loose clothing, and avoiding straining
    \item \textbf{Stool softening:} drinking plenty of water and, where recommended, using fibre supplements or laxatives so that stools pass easily
    \item \textbf{Pain relief:} simple over-the-counter painkillers and, for fissures, topical ointments that relax sphincter spasm
    \item \textbf{Medicines:} treatment of underlying infections or inflammatory conditions as prescribed
    \item \textbf{Procedures:} drainage of abscesses, rubber band ligation for haemorrhoids, and minor surgery for chronic fissures or fistulas
    \item \textbf{Specialist referral:} for persistent or unexplained pain, pelvic floor physiotherapy or review by a colorectal specialist
\end{itemize}

\section*{Complications}
\begin{itemize}
    \item A chronic anal fissure that fails to heal
    \item An abscess that spreads or forms a fistula
    \item Chronic pain or sensitivity affecting daily life and sleep
    \item Anxiety about having a bowel movement, leading to avoidance and worsening constipation
    \item Delayed diagnosis of a more serious underlying condition if the pain is ignored
\end{itemize}

\section*{Prognosis and Outcomes}
Most causes of anal pain settle with simple measures within a few days to weeks. Anal fissures commonly heal within several weeks with stool softening and topical treatment, although a minority become chronic. Haemorrhoid flare-ups usually improve with self-care. Pain caused by muscle spasm may come and go over months but often improves with pelvic floor physiotherapy and relaxation techniques. Outcomes are generally good, especially when the cause is identified early and any serious condition is ruled out.

\section*{Prevention and Self-Care}
\begin{itemize}
    \item Eat a high-fibre diet and drink plenty of fluids to keep stools soft
    \item Do not delay going to the toilet or strain excessively when passing stool
    \item Keep the anal area clean and dry, and avoid harsh soaps or vigorous wiping
    \item Use a cushion or soft seat if sitting is painful
    \item Practise pelvic floor relaxation and manage stress if muscle spasm seems to be a trigger
    \item Seek medical advice early for pain that persists, recurs, or is associated with bleeding or a lump
\end{itemize}

\section*{Sources and Further Reading}
The following resources provide reliable, peer-reviewed and patient-orientated information on anal pain and its causes.

\begin{itemize}
    \item NHS. \textit{Overview of anal pain and common anorectal conditions.} https://www.nhs.uk/conditions/
    \item Mayo Clinic. \textit{Guide to the causes and management of anal pain.} https://www.mayoclinic.org/
    \item MSD Manual. \textit{Professional and patient information on anorectal pain and disorders.} https://www.msdmanuals.com/
    \item MedlinePlus. \textit{Health information on anal pain and anorectal conditions.} https://medlineplus.gov/
    \item NICE. \textit{Clinical guidance on the management of haemorrhoids and anal conditions.} https://www.nice.org.uk/
\end{itemize}
